\documentclass[10pt,a4paper]{amsart}
\usepackage{amsmath,amssymb,amsthm,graphicx}
\usepackage{verbatim}
\newenvironment{reference}{\comment}{\endcomment}
%\newenvironment{reference}{}{}
\newenvironment{slogan}{\comment}{\endcomment}
\newenvironment{history}{\comment}{\endcomment}

% For commutative diagrams we use Xy-pic
\usepackage[all]{xy}

% We use 2cell for 2-commutative diagrams.
\xyoption{2cell}
% Theorem environments.
%
\theoremstyle{plain}
\newtheorem{theorem}[subsection]{Theorem}
\newtheorem{proposition}[subsection]{Proposition}
\newtheorem{lemma}[subsection]{Lemma}

\theoremstyle{definition}
\newtheorem{definition}[subsection]{Definition}
\newtheorem{example}[subsection]{Example}
\newtheorem{exercise}[subsection]{Exercise}
\newtheorem{situation}[subsection]{Situation}

\theoremstyle{remark}
\newtheorem{remark}[subsection]{Remark}
\newtheorem{remarks}[subsection]{Remarks}

\numberwithin{equation}{subsection}

% Macros
%
\def\lim{\mathop{\mathrm{lim}}\nolimits}
\def\colim{\mathop{\mathrm{colim}}\nolimits}
\def\Spec{\mathop{\mathrm{Spec}}}
\def\Hom{\mathop{\mathrm{Hom}}\nolimits}
\def\Ext{\mathop{\mathrm{Ext}}\nolimits}
\def\SheafHom{\mathop{\mathcal{H}\!\mathit{om}}\nolimits}
\def\SheafExt{\mathop{\mathcal{E}\!\mathit{xt}}\nolimits}
\def\Sch{\mathit{Sch}}
\def\Mor{\mathop{\mathrm{Mor}}\nolimits}
\def\Ob{\mathop{\mathrm{Ob}}\nolimits}
\def\Sh{\mathop{\mathit{Sh}}\nolimits}
\def\NL{\mathop{N\!L}\nolimits}
\def\CH{\mathop{\mathrm{CH}}\nolimits}
\def\proetale{{pro\text{-}\acute{e}tale}}
\def\etale{{\acute{e}tale}}
\def\QCoh{\mathit{QCoh}}
\def\Ker{\mathop{\mathrm{Ker}}}
\def\Im{\mathop{\mathrm{Im}}}
\def\Coker{\mathop{\mathrm{Coker}}}
\def\Coim{\mathop{\mathrm{Coim}}}

% Boxtimes
%
\DeclareMathSymbol{\boxtimes}{\mathbin}{AMSa}{"02}

%
% Macros for moduli stacks/spaces
%
\def\QCohstack{\mathcal{QC}\!\mathit{oh}}
\def\Cohstack{\mathcal{C}\!\mathit{oh}}
\def\Spacesstack{\mathcal{S}\!\mathit{paces}}
\def\Quotfunctor{\mathrm{Quot}}
\def\Hilbfunctor{\mathrm{Hilb}}
\def\Curvesstack{\mathcal{C}\!\mathit{urves}}
\def\Polarizedstack{\mathcal{P}\!\mathit{olarized}}
\def\Complexesstack{\mathcal{C}\!\mathit{omplexes}}
% \Pic is the operator that assigns to X its picard group, usage \Pic(X)
% \Picardstack_{X/B} denotes the Picard stack of X over B
% \Picardfunctor_{X/B} denotes the Picard functor of X over B
\def\Pic{\mathop{\mathrm{Pic}}\nolimits}
\def\Picardstack{\mathcal{P}\!\mathit{ic}}
\def\Picardfunctor{\mathrm{Pic}}
\def\Deformationcategory{\mathcal{D}\!\mathit{ef}}
\usepackage{cleveref}
\crefname{originallemma}{original lemma}{original lemmas}
\Crefname{originallemma}{Original lemma}{Original lemmas}
\makeatletter
\newlabel{lemma-quasi-compact-scheme-theoretic-image}{{082Z}{source}{}{sourcepointer-1}{}}
\newlabel{lemma-quasi-compact-scheme-theoretic-image@cref}{{[originallemma][0][]082Z}{[1][0][]source}{}{}{}}
\makeatother
\crefname{originallemma}{original lemma}{original lemmas}
\Crefname{originallemma}{Original lemma}{Original lemmas}
\makeatletter
\newlabel{lemma-scheme-theoretic-image-reduced}{{0830}{source}{}{sourcepointer-2}{}}
\newlabel{lemma-scheme-theoretic-image-reduced@cref}{{[originallemma][0][]0830}{[1][0][]source}{}{}{}}
\makeatother
\crefname{originallemma}{original lemma}{original lemmas}
\Crefname{originallemma}{Original lemma}{Original lemmas}
\makeatletter
\newlabel{morphisms-lemma-quasi-compact-immersion}{{01RG}{source}{}{sourcepointer-3}{}}
\newlabel{morphisms-lemma-quasi-compact-immersion@cref}{{[originallemma][0][]01RG}{[1][0][]source}{}{}{}}
\makeatother
\title{An immersion with quasi-compact source or reduced source is an open immersion onto a dense open of its scheme theoretic image}
\author{The Stacks Project Authors}
\date{stacks-project:spaces-morphisms.tex:local-snapshot-20261012; GNU FDL 1.2 or later}
\begin{document}
\maketitle
\setcounter{subsection}{1}
\noindent\textit{Editorial note.} All delivered mathematical text is the author's own, copied from the pinned Stacks Project chapter file spaces-morphisms.tex; no statement or proof step was rewritten, reordered or removed.\par\medskip
\noindent\textit{Editorial note.} The statement carries the disjunctive hypothesis (quasi-compact or reduced) and an 'Moreover' clause on topological density; both are inside the counted statement span.\par\medskip
\noindent\textit{Editorial note.} The three references of the counted spans are recorded as non-load-bearing navigation pointers; the argument that consumes them stays inside the counted proof span.\par\medskip
\noindent\textit{Editorial note.} Preamble: the three native excerpt spans of sources/stacks-spaces-morphisms/preamble.tex (lines 8-13, 14-18 and 154-217). The chapter file itself carries no preamble text.\par\medskip
\noindent\textit{Allowed standard prerequisites (editorial).}\begin{itemize}
\item Commutation of the scheme theoretic image with etale localization in the quasi-compact case. Lemma lemma-quasi-compact-scheme-theoretic-image (3249-3265), used at proof lines 3580-3582. Handles the first alternative of the hypothesis.
\item The scheme theoretic image of a morphism from a reduced space. Lemma lemma-scheme-theoretic-image-reduced (3286-3296), used at proof lines 3582-3583. Handles the second alternative of the hypothesis.
\item The scheme case for a quasi-compact immersion. Morphisms, Lemma morphisms-lemma-quasi-compact-immersion, used at proof lines 3584-3585. Is the input that closes the argument.
\end{itemize}\medskip
\ifdefined\hypertarget\hypertarget{sourcepointer-1}{}\fi
\noindent\textit{Source reference.} Original lemma 082Z: focused-seeds/AGP/sources/stacks-snapshot-20261011/tags.txt, source lines 10269--10269. Navigation pointer into the Stacks Project; the standard fact it supplies is applied inside the counted proof span and the argument resting on it stays there.\par\medskip
\ifdefined\hypertarget\hypertarget{sourcepointer-2}{}\fi
\noindent\textit{Source reference.} Original lemma 0830: focused-seeds/AGP/sources/stacks-snapshot-20261011/tags.txt, source lines 10270--10270. Navigation pointer into the Stacks Project; the standard fact it supplies is applied inside the counted proof span and the argument resting on it stays there.\par\medskip
\ifdefined\hypertarget\hypertarget{sourcepointer-3}{}\fi
\noindent\textit{Source reference.} Original lemma 01RG: focused-seeds/AGP/sources/stacks-snapshot-20261011/tags.txt, source lines 2301--2301. Navigation pointer into the Stacks Project; the standard fact it supplies is applied inside the counted proof span and the argument resting on it stays there.\par\medskip
\begin{lemma}
\label{lemma-quasi-compact-immersion}
Let $S$ be a scheme. Let $h : Z \to X$ be an immersion of algebraic spaces
over $S$. Assume either $Z \to X$ is quasi-compact or $Z$ is reduced.
Let $\overline{Z} \subset X$ be the scheme theoretic image of $h$.
Then the morphism $Z \to \overline{Z}$ is an open immersion
which identifies $Z$ with a scheme theoretically dense open
subspace of $\overline{Z}$. Moreover, $Z$ is topologically
dense in $\overline{Z}$.
\end{lemma}

\begin{proof}
In both cases the formation of $\overline{Z}$ commutes with
\'etale localization, see
Lemmas \ref{lemma-quasi-compact-scheme-theoretic-image} and
\ref{lemma-scheme-theoretic-image-reduced}.
Hence this lemma follows from the case of schemes, see
Morphisms, Lemma \ref{morphisms-lemma-quasi-compact-immersion}.
\end{proof}

\end{document}
